% =====================================================================
\chapter{Evaluación: entropía cruzada y perplejidad}
\label{tema:5}
% =====================================================================

Un buen modelo asigna probabilidad alta al texto real. Este tema define la pérdida que se minimiza al entrenar todos los modelos de lenguaje, la entropía cruzada, y la medida de evaluación derivada de ella, la perplejidad. Se sigue~\cite[§3.2, §3.3 y §3.7]{jm3}; los resultados de teoría de la información proceden de~\cite[cap.~2]{ct}.

\section{Pérdida y perplejidad}

Sea un texto de evaluación con $N$ predicciones y sea $p_t=P(x_t\mid x_{1:t-1})\in(0,1]$ la probabilidad que el modelo asigna al token que realmente aparece en la posición $t$.

\begin{definicion}[Pérdida y perplejidad]\label{def:ppl}
La \emph{pérdida} (log-verosimilitud negativa media) y la \emph{perplejidad} del modelo sobre el texto son
\[
\mathcal{J}=-\frac{1}{N}\sum_{t=1}^{N}\ln p_t,\qquad \PPL=e^{\mathcal{J}}.
\]
Con logaritmo natural, $\mathcal{J}$ se mide en \emph{nats}; con $\log_2$, en \emph{bits}, y $\mathcal{J}_{\text{bits}}=\mathcal{J}_{\text{nats}}/\ln 2$.
\end{definicion}

\begin{proposicion}\label{prop:ppl}
\begin{enumerate}
  \item $\PPL=\bigl(\prod_{t=1}^{N}1/p_t\bigr)^{1/N}$, la media geométrica de los inversos de las probabilidades.
  \item $\PPL\ge 1$, con igualdad si y solo si $p_t=1$ para todo $t$.
  \item Un modelo uniforme sobre $K$ tokens tiene $\PPL=K$ sobre cualquier texto.
\end{enumerate}
\end{proposicion}

\begin{proof}
(i) $-\ln p_t=\ln(1/p_t)$, y la suma de logaritmos es el logaritmo del producto: $\mathcal{J}=\frac1N\ln\prod_t(1/p_t)=\ln\bigl(\prod_t 1/p_t\bigr)^{1/N}$; basta aplicar la exponencial. (ii) Cada $1/p_t\ge 1$; el producto y su raíz $N$-ésima son $\ge 1$, con igualdad si y solo si todos los factores valen~1. (iii) Si $p_t=1/K$ para todo $t$, el producto es $K^{N}$ y su raíz $N$-ésima es $K$.
\end{proof}

La perplejidad admite una interpretación como \emph{factor de ramificación medio}: un modelo con perplejidad $k$ está, en promedio, tan indeciso como si eligiera uniformemente entre $k$ tokens~\cite[§3.3.1]{jm3}.

\begin{ejemplo}\label{ej:ppl}
El modelo del \cref{ej:bigramas}, evaluado sobre su propio corpus (13 bigramas), tiene $\sum_t\ln p_t=-8{,}047$. Por tanto, $\mathcal{J}=8{,}047/13=0{,}619$ nats $=0{,}893$ bits y $\PPL=e^{0{,}619}=1{,}86$. El modelo uniforme sobre los 5 tokens tiene $\mathcal{J}=\ln 5=1{,}609$ nats y $\PPL=5$.
\end{ejemplo}

\section{Entropía y entropía cruzada}

\begin{definicion}[Entropía y entropía cruzada]
Sean $p$ y $q$ distribuciones sobre un conjunto finito $\mathcal{X}$, con $q(x)>0$ siempre que $p(x)>0$. La \emph{entropía} de $p$ y la \emph{entropía cruzada} de $q$ respecto de $p$ son
\[
H(p)=-\sum_{x\in\mathcal{X}}p(x)\ln p(x),\qquad
H(p,q)=-\sum_{x\in\mathcal{X}}p(x)\ln q(x),
\]
con el convenio $0\ln 0=0$.
\end{definicion}

Si los tokens del texto se distribuyen según $p$ y el modelo es $q$, la pérdida $\mathcal{J}$ es un estimador de $H(p,q)$: es la media empírica de $-\ln q(X)$ con $X\sim p$.

\begin{teorema}[Desigualdad de Gibbs]\label{teo:gibbs}
$H(p,q)\ge H(p)$, con igualdad si y solo si $q=p$.
\end{teorema}

\begin{proof}
Sea $S=\{x:p(x)>0\}$. Entonces
\[
H(p)-H(p,q)=\sum_{x\in S}p(x)\ln\frac{q(x)}{p(x)}\le\ln\sum_{x\in S}p(x)\frac{q(x)}{p(x)}=\ln\sum_{x\in S}q(x)\le\ln 1=0,
\]
donde la primera desigualdad es la de Jensen para la función estrictamente cóncava $\ln$. La igualdad en Jensen exige que $q(x)/p(x)$ sea constante en $S$, y la igualdad final exige $\sum_{x\in S}q(x)=1$; ambas juntas dan $q=p$~\cite[teorema~2.6.3]{ct}.
\end{proof}

La diferencia $H(p,q)-H(p)=\sum_x p(x)\ln\bigl(p(x)/q(x)\bigr)$ es la \emph{divergencia de Kullback–Leibler} $D(p\,\|\,q)$, que se estudia en el módulo~6.

\begin{corolario}\label{cor:aprender}
Entre todos los modelos $q$, la entropía cruzada esperada $H(p,q)$ se minimiza exactamente en $q=p$, y su valor mínimo es $H(p)$.
\end{corolario}

El \cref{cor:aprender} justifica el criterio de entrenamiento de los modelos de lenguaje: minimizar la entropía cruzada media sobre un corpus empuja al modelo a reproducir la distribución real del token siguiente. Ningún modelo puede bajar de $H(p)$, la incertidumbre propia del lenguaje. La relación entre entropía y predicción se remonta a Shannon~\cite{shannon}.

\section{Entrenamiento y evaluación}

\begin{observacion}
La perplejidad debe medirse sobre un conjunto de evaluación disjunto del de entrenamiento. Evaluar sobre los datos de entrenamiento da una estimación optimista: el EMV del \cref{teo:emv} maximiza precisamente la verosimilitud de esos datos~\cite[§3.2]{jm3}. Además, la perplejidad solo es comparable entre modelos que usan el mismo vocabulario, porque depende de la tokenización.
\end{observacion}

\begin{lecturas}
\item Jurafsky y Martin~\cite{jm3}, §3.2 «Evaluating Language Models: Training and Test Sets», §3.3 «Evaluating Language Models: Perplexity» y §3.7 «Perplexity's Relation to Entropy».
\item Cover y Thomas~\cite{ct}, §2.1 «Entropy», §2.3 «Relative Entropy and Mutual Information» y §2.6 «Jensen's Inequality and Its Consequences».
\item MacKay~\cite{mackay}, cap.~2 (entropía y desigualdad de Gibbs).
\item Zhang \emph{et al.}~\cite{d2l}, §9.3.2 «Perplexity».
\end{lecturas}

\section*{Ejercicios del tema 5}
\addcontentsline{toc}{section}{Ejercicios del tema 5}

\begin{ejercicio}{1}{eje:5-ppl1}
¿Qué significa que un modelo tenga perplejidad~1 sobre un texto? Justifíquelo con la \cref{prop:ppl}.
\end{ejercicio}

\begin{ejercicio}{2}{eje:5-perdida}
Un modelo asigna las probabilidades $1/2$, $1/4$, $1/4$ y $1/8$ a los cuatro tokens de un texto de evaluación. Calcule la pérdida en bits y en nats, y la perplejidad.
\end{ejercicio}

\begin{ejercicio}{2}{eje:5-digitos}
Un conjunto de entrenamiento contiene 100 dígitos: 82 ceros y dos de cada dígito del 1 al 9. Calcule la perplejidad de un modelo de unigramas de máxima verosimilitud sobre el texto de prueba «0 0 0 0 0 3 0 0 0 0» y compárela con la de un modelo uniforme sobre los diez dígitos.
\fuente{adaptado de Jurafsky y Martin~\cite{jm3}, cap.~3, ejercicio~3.12 (conteos modificados).}
\end{ejercicio}

\begin{ejercicio}{3}{eje:5-ppl}
Demuestre la \cref{prop:ppl} con todo detalle, justificando cada propiedad del logaritmo que utilice.
\end{ejercicio}

\begin{ejercicio}{3}{eje:5-lagrange}
Demuestre el \cref{cor:aprender} sin usar la desigualdad de Jensen, mediante multiplicadores de Lagrange: minimice $H(p,q)$ en $q$ con la restricción $\sum_x q(x)=1$, suponiendo $p(x)>0$ para todo $x$.
\end{ejercicio}
